% !TeX root = . . /main. tex

\chapter{中学常见恒等式及一般证法}
\section{中学常见数学恒等式类型}
中学数学中恒等式的类型并不繁杂. 常见的恒等式类型有以下三种:

代数恒等式: 代数恒等式指在变量取值范围内, 等号两边始终相等的等式. 代数恒等式通常是可通过代数变形验证的. 代数恒等式常见类型有形如对数恒等式、指数恒等式、完全平方公式等类型. 

% 三角恒等式: 指在三角函数定义域内, 对定义域内所有取值都成立的等式. 通常用于描述三角函数之间的联系或者是简化三角函数表达式. 三角恒等式有诱导公式、 和差化积公式等类型. 
% 三角恒等式: 是指在定义域内, 对定义域内所有取值都成立的恒等式等式. 通常是用来描述三角函数之间的联系或者是简化三角函数表达式. 常见的类型有诱导公式、和差化积公式等类型. 
三角恒等式: 指在定义域内, 对定义域内所有取值都成立的恒等式. 通常是用来描述三角函数不同符号间的联系或者是简化三角函数表达式. 常见类型有诱导公式、和差化积公式等类型. 


组合恒等式: 是指含有到组合数或者组合结构的恒等式. 通常反映了组合数的内在性质或者是不同结构的等价关系. 常见类型有求和恒等式、组合数对称性恒等式等类型. 


\section{恒等式证明的传统方法概述}
在中学数学中恒等式证明的传统方法主要有三种一一代数法、几何法以及数学归纳法. 

代数法是最日常最容易被掌握的方法. 它的主要依据是代数式的基本运算规则和性质, 通过对等式两边同时进行加、减、乘、除等运算, 或者通过因式分解、移项等基本操作, 逐步操作使等式两边的形式一致, 从而完成证明. 

% 几何法是恒等式证明方法中最直观的方法. 通过将恒等式中的元素转化为几何量, 通过对几何图形的构造推理, 实现恒等式的证明, 从而证明恒等式. 但相对代数法和数学归纳法来说, 几何法的局限性比较大. 
几何法是恒等式证明方法中最直观的方法. 原理是将恒等式中的元素转化为几何量, 构造一个适合恒等式结构的几何图形, 复现恒等式, 从而证明这个恒等式. 相对代数法和数学归纳法来说, 几何法的局限性比较大，适用范围更小. 


数学归纳法是恒等式证明方法中最严谨的方法. 先验证当 \(n\) 取第一个值 \({n}_{0}\) (通常为 1 )时恒等式成立,确保了起点的正确性. 然后假设当 \(n = k\) ( \(k\) 为自然数且 \(k \geq  {n}_{0}\) )时恒等式成立,在此基础上,通过计算 \(n = k + 1\) 时的结果,恒等式成立就完成了证明. 

\section{恒等式证明的传统方法范例}

\begin{example}\label{example:例1}
  证明恒等式: \[1 + \frac{A - a}{A - 1} + \frac{\left( {A - a}\right) \left( {A - a - 1}\right) }{\left( {A - 1}\right) \left( {A - 2}\right) } + \ldots  + \frac{\left( {A - a}\right) \cdots 2 \cdot  1}{\left( {A - 1}\right) \cdots \left( {a + 1}\right) a} = \frac{A}{a} ,\]
  其中 \(A,a\) 都是正整数,且 \(A > a\) . 
\end{example}

% \begin{proof}
%   恒等式左右两边同乘 \(\left( {A - 1}\right) \left( {A - 2}\right) \cdots \left( {a + 1}\right) a\) 得,

% 等式左边 \(=\)

% \(\left( {A - 1}\right) \left( {A - 2}\right) \cdots \left( {a + 1}\right) a + \left( {A - a}\right) \left( {A - a - 1}\right) \left( {A - 2}\right) \cdots \left( {a + 1}\right) a + \ldots  + \left( {A - a}\right) \left( {A - a - 1}\right) \cdots 2 \cdot  1\)

% 等式右边 \(= A\left( {A - 1}\right) \left( {A - 2}\right) \left( {A - 3}\right) \cdots \left( {a + 1}\right)\)

% 左边后两项相加 \(= a\left( {A - a}\right) \left( {A - a - 1}\right) \cdots 2 + \left( {A - a}\right) \left( {A - a - 1}\right) \cdots 2 \cdot  1\)

% \[
% = \left( {a + 1}\right) \left( {A - a}\right) \left( {A - a - 1}\right) \cdots 2
% \]

% 左边后三项相加 \(= 2\left( {a + 1}\right) \left( {A - a}\right) \left( {A - a - 1}\right) \cdots 3 + a\left( {a + 1}\right) \left( {A - a}\right) \left( {A - a - 1}\right) \cdots 3\)

% \[
% = \left( {a + 1}\right) \left( {a + 2}\right) \left( {A - a}\right) \left( {A - a - 1}\right) \cdots 3
% \]

% 以此类推,可得最后等式左边 \(= \left( {a + 1}\right) \left( {a + 2}\right) \cdots \left( {A - 2}\right) \left( {A - 1}\right) A\)

% 等式左右两边相等. 
% \end{proof}

\begin{proof}
  恒等式左右两边同乘 \(\left( {A - 1}\right) \left( {A - 2}\right) \cdots \left( {a + 1}\right) a\) 得：

  等式左边 
  $
  \begin{aligned}[t]
      =&\left( {A - 1}\right) \left( {A - 2}\right) \cdots \left( {a + 1}\right) a \\
      &+ \left( {A - a}\right) \left( {A - a - 1}\right) \left( {A - 2}\right) \cdots \left( {a + 1}\right) a \\
      &+ \cdots \\
      &+ \left( {A - a}\right) \left( {A - a - 1}\right) \cdots 2 \cdot 1
  \end{aligned}$

  等式右边 \(= A\left( {A - 1}\right) \left( {A - 2}\right) \left( {A - 3}\right) \cdots \left( {a + 1}\right)\)

  左边后两项相加：
  \begin{align*}
      &a\left( {A - a}\right) \left( {A - a - 1}\right) \cdots 2 + \left( {A - a}\right) \left( {A - a - 1}\right) \cdots 2 \cdot 1 \\
      &= \left( {a + 1}\right) \left( {A - a}\right) \left( {A - a - 1}\right) \cdots 2
  \end{align*}

  左边后三项相加：
  \begin{align*}
      &2\left( {a + 1}\right) \left( {A - a}\right) \left( {A - a - 1}\right) \cdots 3 + a\left( {a + 1}\right) \left( {A - a}\right) \left( {A - a - 1}\right) \cdots 3 \\
      &= \left( {a + 1}\right) \left( {a + 2}\right) \left( {A - a}\right) \left( {A - a - 1}\right) \cdots 3
  \end{align*}

  以此类推，可得最后等式左边 \(= \left( {a + 1}\right) \left( {a + 2}\right) \cdots \left( {A - 2}\right) \left( {A - 1}\right) A\)

  等式左右两边相等。
\end{proof}



\begin{example}
  几何法证明 \({\left( a + b\right) }^{2} = {a}^{2} + {2ab} + {b}^{2}\) . 
\end{example}

\begin{proof}
  构造一个边长为 \(\left( {a + b}\right)\) 的大正方形. 其面积就为 \(S = {\left( a + b\right) }^{2}\) . 将正方形相邻两边分别进行分割,分割为 \(a,b\) 两段. 大正方形的面积就分割成了四个部分分别为 \({a}^{2}\) , \({b}^{2},{ab},{ab}\) . 面积相等可知 \({\left( a + b\right) }^{2} = {a}^{2} + {b}^{2} + {ab} + {ab}\), 即 ${\left( a + b\right) }^{2} = {a}^{2} + {2ab} + {b}^{2}$. 
\end{proof}

\begin{example}
  数学归纳法证明恒等式 \((a + b)^n = \sum_{k=0}^n C_n^k a^{n-k}b^k\)，\(a, b, n\) 均为自然数。
\end{example}

\begin{proof}
当 \(n = 1\) 时，\((a + b)^1 = C_{1}^{0} a^{1} b^{0} + C_{1}^{1} a^{0} b^{1} = a + b\)，等式成立. 

假设当 \(n = m\) 时，等式成立，即有 \((a + b)^m = \sum_{k = 0}^{m} C_{m}^{k} a^{m - k} b^{k}\). 

当 \(n = m + 1\) 时，
\begin{align*}
(a + b)^{m + 1} &= (a + b) \cdot (a + b)^m \\
&= (a + b) \left( \sum_{k = 0}^{m} C_{m}^{k} a^{m - k} b^{k} \right) \\
&= \sum_{k = 0}^{m} C_{m}^{k} a^{m - k + 1} b^{k} + \sum_{k = 0}^{m} C_{m}^{k} a^{m - k} b^{k + 1} \\
&= a^{m + 1} + \sum_{k = 1}^{m} C_{m}^{k} a^{m - k + 1} b^{k} + \sum_{k = 1}^{m} C_{m}^{k - 1} a^{m - k + 1} b^{k} + b^{m + 1} \\
&= \sum_{k = 0}^{m + 1} \left( C_{m}^{k} + C_{m}^{k - 1} \right) a^{(m + 1) - k} b^{k} \\
&= \sum_{k = 0}^{m + 1} C_{m + 1}^{k} a^{(m + 1) k} b^{k}. 
\end{align*}
等式成立。因此，当 \(a,b,n\) 是自然数时，恒等式 \((a + b)^n = \sum_{k = 0}^{n} C_{n}^{k} a^{n - k} b^{k}\) 成立. 
\end{proof}

% \begin{example}
%   数学归纳法证明恒等式 \({1}^{2} + {2}^{2} + {3}^{2} + \ldots  + {n}^{2} = \frac{n\left( {n + 1}\right) \left( {{2n} + 1}\right) }{6}\) . 
% \end{example}

% \begin{proof}
%   \(n = 1\) 时，等式左边 \(= 1\)，等式右边 \(= 1\)，等式成立。

%   假设 \(n = k\)（\(k\) 是正整数）时等号成立，就有
%   \[
%   1^{2} + 2^{2} + 3^{2} + \cdots + k^{2} = \frac{k(k + 1)(2k + 1)}{6}
%   \]
%   成立。当 \(n = k + 1\) 时，就有
%   \begin{align*}
%       1^{2} + 2^{2} + 3^{2} + \cdots + k^{2} + (k + 1)^{2} 
%       &= (1^{2} + 2^{2} + 3^{2} + \cdots + k^{2}) + (k + 1)^{2} \\
%       &= \frac{k(k + 1)(2k + 1)}{6} + (k + 1)^{2} \\
%       &= (k + 1)\left[\frac{k(2k + 1)}{6} + k + 1\right] \\
%       &= (k + 1)\left(\frac{2k^{2} + k + 6k + 6}{6}\right) \\
%       &= (k + 1)\left(\frac{2k^{2} + 7k + 6}{6}\right) \\
%       &= (k + 1)\frac{(k + 2)(2k + 3)}{6} \\
%       &= \frac{(k + 1)[(k + 1) + 1][2(k + 1) + 1]}{6}
%   \end{align*}
%   假设成立，恒等式成立。
% \end{proof}
% \begin{proof}
%   \(n = 1\) 时,等式左边 \(= 1\) ,等式右边 \(= 1\) ,等式成立. 

% 假设 \(n = k\) ( \(k\) 是正整数) 时等号成立, 就有
% \[
% {1}^{2} + {2}^{2} + {3}^{2} + \ldots  + {k}^{2} = \frac{k\left( {k + 1}\right) \left( {{2k} + 1}\right) }{6}
% \]
% 成立.  当 \(n = k + 1\) 时,就有
% \[
% {1}^{2} + {2}^{2} + {3}^{2} + \ldots  + {k}^{2} + {\left( k + 1\right) }^{2}
% \]

% \[
% = \left( {{1}^{2} + {2}^{2} + {3}^{2} + \ldots  + {k}^{2}}\right)  + {\left( k + 1\right) }^{2}
% \]

% \[
% = \frac{k\left( {k + 1}\right) \left( {{2k} + 1}\right) }{6} + {\left( k + 1\right) }^{2}
% \]

% \[
% = \left( {k + 1}\right) \left\lbrack  {\frac{k\left( {{2k} + 1}\right) }{6} + k + 1}\right\rbrack
% \]

% \[
% = \left( {k + 1}\right) \left( \frac{2{k}^{2} + k + {6k} + 6}{6}\right)
% \]

% \[
% = \left( {k + 1}\right) \left( \frac{2{k}^{2} + {7k} + 6}{6}\right)
% \]

% \[
% = \left( {k + 1}\right) \frac{\left( {k + 2}\right) \left( {{2k} + 3}\right) }{6}
% \]

% \[
% = \frac{\left( {k + 1}\right) \left\lbrack  {\left( {k + 1}\right)  + 1}\right\rbrack  \left\lbrack  {2\left( {k + 1}\right)  + 1}\right\rbrack  }{6}
% \]

% 假设成立. 恒等式成立. 
% \end{proof}

在这里的三个例题, 分别选用了三个方法的例题各一个, 是为了展示三个方法各自的基本流程. 也是为了和后面概率模型法做一个简单的对比. 
